\section{Determinant ideals and spherical heights}
\label{sec:chow-metric}
We prove the multiform estimates and the norm-to-degree step in
\eqref{eq:chow-slope-raw}. Let \(F\) be a number field, \(e_F=[F:\mathbb Q]\),
and \(w:F^n\twoheadrightarrow F^N\) a spanning coordinate matrix with
\(1\le N\le n\). At each embedding \(\sigma:F\hookrightarrow\mathbb C\)
let \(A_\sigma\) be an invertible frame, so the norm of \(x\in\mathbb C^N\)
in the coordinate space is \(\|A_\sigma^{-1}x\|_2\).
Assume every column of \(\sigma w\) has frame norm at most
\(\Lambda\ge1\). Form heights will use the mean over \emph{all} \(e_F\)
embeddings, counting a conjugate pair twice.

For a nonzero polynomial \(g\) with \(B\) blocks of \(N\) variables let
\(J(g)\) be the fractional ideal generated by the coefficients of the
polynomial obtained by substituting \(w\) in every block. The substituted
form is nonzero because \(w\) has a right inverse. Write \(m_\sigma(g)\)
for the expectation of \(\log|(\sigma g)(A_\sigma z_1,\ldots,A_\sigma z_B)|\),
where the \(z_i\) are independent uniform points on the unit sphere
in \(\mathbb C^N\), of real dimension \(2N-1\). Set
\[
 H(g)=-\frac{\log \mathrm N J(g)}{e_F}
               +\frac1{e_F}\sum_\sigma m_\sigma(g).
\]
This uses the norm on fractional ideals extended multiplicatively.
For a multihomogeneous form, its nonzero restriction to the product of
spheres has logarithm locally integrable and its zero set has measure
zero. One may use logarithmic integrability of a nonzero holomorphic
polynomial in complex affine charts: homogeneity reduces the statement
to the compact projective directions, and the radial and phase factors
are smooth in every unit chart. All expectations below are consequently
well defined. At a finite prime the minimum coefficient valuation is
additive on products: after dividing each factor by its Gauss norm
in a local uniformizer, the reductions have nonzero product over the
residue field. Hence \(J(gh)=J(g)J(h)\). At infinity the logarithmic
means add. The product formula gives \(H(cg)=H(g)\) for \(c\in F^\times\).

\paragraph{An evaluation bound including its constant.}
Let \(g\) have block degrees \(a_1,\ldots,a_B\), and put
\(a=\sum_i a_i\), \(H_{N-1}=\sum_{k=1}^{N-1}1/k\). For unit vectors
\(v_i\in\mathbb C^N\), unitary invariance and Jensen's inequality on
a circle give
\[
 \log|(\sigma g)(A_\sigma v_1,\ldots,A_\sigma v_B)|
      \le m_\sigma(g)+a H_{N-1}/2
\]
whenever the value on the left is nonzero. To see the constant, rotate
each \(v_i\) to the first coordinate, changing the polynomial by the
inverse unitary change of variables. In one block write a random vector
as \(z=(z_1,z')\). For fixed \(z_1,z'\), average
\(\log|g(z_1,e^{it}z')|\) over \(t\). The submean inequality bounds
this average below by \(\log|g(z_1,0)|\), allowing \(-\infty\) for
a vanishing evaluation. Repeat over all blocks. The left average
integrated over the spheres is unchanged by rotation, while the final
value equals the chosen evaluation times
\(\prod_i z_{i1}^{a_i}\). For \(N>1\), \(|z_1|^2\) has density
\((N-1)(1-u)^{N-2}\) on \(0<u<1\); this follows by normalizing
\(N\) independent standard complex Gaussians, whose squared radii
are independent exponentials. Its logarithmic loss is
\[
 \mathbb E[-\log |z_1|^2]
    =\int_0^1\frac{1-(1-t)^{N-1}}t\,dt
    =\sum_{j=0}^{N-2}\int_0^1(1-t)^j\,dt=H_{N-1}.
\]
For \(N=1\) the loss is zero. This proves the claim.

On the unit polydisk of the \(n\) substitution variables in each block,
the substitution \(\sigma w\) produces vectors of frame norm at most
\(n\Lambda\), by the triangle inequality. Homogeneity and the bound just
proved give \(|(\sigma g)(\sigma w z)|\le
\exp(m_\sigma(g)+a H_{N-1}/2)(n\Lambda)^a\).
By the coefficient formula on the unit torus the same bound holds for
each coefficient of \((\sigma g)(\sigma w z)\).
Take any nonzero coefficient \(\alpha\) of \(g(wz)\).
The inclusion \((\alpha)\subset J(g)\) gives
\(\log \mathrm N J(g)\le\sum_\sigma\log|\sigma\alpha|\).
Using the coefficient bounds and \(H_{N-1}\le1+\log n\) yields
\[
 H(g)\ \ge\ -a\left(\log(n\Lambda)+\frac{\log n+1}{2}\right).
\]
When \(g=f\) is a Chow form with \(D+1\) blocks, all of degree \(\delta\),
the conventional form height in these frames adds
\((D+1)\delta H_{N-1}/2\) to \(H(f)\).
Dividing by \((D+1)\delta\) gives precisely
\(h_f=h-H_{N-1}/2\ge h-(1+\log n)/2\), with the height \(h\) of the
image. Replacing its defining form by a scalar multiple or a positive
power leaves the normalized height unchanged when the degree is adjusted.

\paragraph{Lattice degree from polynomial columns.}
Let \(M\) be a full \(\mathcal O_F\)-lattice in a \(d\)-dimensional
\(F\)-space, \(b\) an \(F\)-basis and \(G_\sigma\) its positive Hermitian
Gram matrix. Define \(I_b(M)\) as the fractional ideal spanned by the
determinants of all \(d\)-tuples of lattice elements relative to \(b\).
With this convention
\[
 \widehat{\deg}(M)=-\frac{\log\mathrm N I_b(M)}{e_F}
                 -\frac1{2e_F}\sum_\sigma\log\det G_\sigma .
\]
Indeed, after localization at a finite prime the lattice is free, and
the local determinant ideal records the sum of its elementary-divisor
valuations relative to the basis lattice. Their sum weighted by residue
field logarithms gives the finite term. At infinity the norm of
\(b_1\wedge\cdots\wedge b_d\) is \(\sqrt{\det G_\sigma}\).
Under a change of basis by \(S\), the ideal is divided by \(\det S\)
and the Gram determinant is multiplied by \(|\sigma\det S|^2\);
the product formula cancels the two changes. This also covers a
projective lattice with nonprincipal determinant ideal.

Let \(P\) be a square matrix of polynomials, whose columns are vectors
in the \(b\)-coordinate space, and set \(g=\det P\ne0\).
After substituting \(w\), suppose every coefficient vector of every
column lies in \(M\). Expanding by multilinearity, the coefficient
of each monomial in the determinant is a sum of determinants of such
coefficient vectors. Hence \(J(g)\subset I_b(M)\).
At a complex place suppose column \(k\) has Hermitian norm at most
\(B_k>0\) on unit input vectors in every frame block. Hadamard's
inequality in the Hermitian inner product gives pointwise
\[
 |\det P|\,\sqrt{\det G_\sigma}\le\prod_{k=1}^{d}B_k.
\]
Taking the logarithm almost everywhere and integrating gives
\(m_\sigma(g)\le\sum_k\log B_k-\tfrac12\log\det G_\sigma\).
Since \(J(g)\subset I_b(M)\) implies
\(-\log\mathrm N J(g)\le-\log\mathrm N I_b(M)\),
summing these inequalities proves
\[
 H(\det P)-\sum_k\log B_k\le\widehat{\deg}(M).
\]
Taking the three families of columns and their bounds from
Section~\ref{sec:chow-determinant} proves \eqref{eq:chow-slope-raw},
including its finite-place contribution.

\paragraph{Source correspondences.}
The lower bound and harmonic correction are
\srcid{MultiformHeightLowerBound} and \srcid{FormHeightEqHt}.
The determinant and its norms are
\srcid{LatticeDegreeGeHeightPolynomialVectorsDet},
\srcid{CoeffDetMemSpanDetIntegralPolynomialVectors}, and
\srcid{FiniteHeightSubMeanLeDegreeFullLatticeDetIdeal}.
